\documentclass[10pt,letterpaper]{article}
\usepackage{cornell-notes}

\title{Clause 26.07.09.03: Class Template Transform View Iterator}
\author{}
\date{}
\setDocTitle{Clause 26.07.09.03: Class Template Transform View Iterator}
\setDocSubtitle{C++ 2024}
\setDocOwner{C++ 2024 Cornell Notes}
\setDocDate{\today}
\setCornellCollection{C++ 2024 Cornell Notes}
\setCornellUnitType{Clause}
\setCornellUnitNumber{26.07.09.03}
\setCornellUnitTitle{Class Template Transform View Iterator}
\hypersetup{pdftitle={Clause 26.07.09.03: Class Template Transform View Iterator},pdfsubject={Cornell notes},pdfauthor={}}

\begin{document}
\maketitle
\begin{CornellOverviewBox}{Overview}
\textbf{Classification:} Standard library - Normative\\[2pt]
\textbf{Learning objective:} Understand the rules, terminology, and semantic consequences associated with Class template transform\_view::iterator.
\end{CornellOverviewBox}\begin{CornellSummaryBox}{Summary}
{\sffamily\bfseries\color{Accent} Summary}\par\vspace{3pt}Section 26.7.9.3 focuses on Class template transform\_view::iterator. Apply the specified interface together with its mandates, effects, result conditions, exception behavior, and complexity constraints. When applying it, preserve the distinction between mandatory requirements, recommendations, permissions, and behavior deliberately left undefined, unspecified, or implementation-defined.
\end{CornellSummaryBox}

\end{document}
